\section{Conclusions, Limitations and Future Work}
\label{sec:conclusions}

\subsection{Conclusions}

We presented a controlled, iterative study of the EPIC-KITCHENS-100 Multi-Instance Retrieval challenge in which three progressively stronger approaches were evaluated under a single Codabench protocol. Starting from a JPoSE feature-based baseline at $53.53$\,nDCG, a score-level ensemble of three frozen retrieval models combined with graph-diffusion re-ranking lifts performance to $55.82$\,nDCG, while an AVION ViT-L/14 dual encoder fine-tuned with the SMS loss and complemented by horizontal-flip plus 32-frame test-time augmentation reaches $69.68$\,nDCG over sixteen epochs of single-GPU training. The final iteration recovers more than half of the gap between feature-based baselines and the current public state of the art (CR-CLIP, $82.10$\,nDCG~\cite{he2025crclip}) using a single A100 instance and the publicly released SMS-Loss training code. The per-iteration breakdown isolates where each gain originates. Encoder capacity coupled with relevance-aware supervision dominates by an order of magnitude, while inference-time strategies contribute small but reliable deltas at zero training cost. Additionally, this provides a transparent compute-allocation guideline for future EK-100 MIR work operating under similar resource constraints.

\subsection{Limitations}
\label{sec:conclusions:limitations}

Two limitations frame our results. First, all Approach~3 experiments were performed on a single NVIDIA A100 (40\,GB) with a batch size of $48$, which made multi-GPU runs, larger effective batch sizes and multi-seed re-training infeasible within the project timeline. Several knobs that are routine in the EK-100 MIR literature were therefore left unexplored. These include the contribution of larger batches (which directly affects the SMS gradient through the implicit hard-negative mining in Eq.~\eqref{eq:sms}), systematic sweeps over $\theta$, $\tau$, $\beta_p, \beta_n$, longer fine-tuning schedules above 16 epochs, and variance estimation through repeated runs with different random seeds. Second, the EK-100 challenge releases the soft-label relevance matrix only for the training split. As mentioned previously, the test relevance matrix required to compute nDCG and mAP offline is held out by the organizers and is accessible exclusively through the Codabench grader. Model selection during fine-tuning therefore had to rely on training-loss trajectories rather than on a relevance-aligned validation signal, and every reported number in Table~\ref{tab:iterations} required a Codabench submission, which capped the practical iteration budget. As a sanity check on the JPoSE baseline implementation we did obtain offline numbers on the train-set validation routine ($0.690$\,nDCG, $0.734$\,mAP), but these are not directly comparable to the Codabench score for the same submission ($53.53$\,nDCG) and were not used for cross-iteration comparison.

\subsection{Future work}
\label{sec:conclusions:future}

The most promising extension consistent both with the empirical evidence in Table~\ref{tab:iterations} and with the SMS-Loss recipe of Wang et al.~\cite{wang2024symmetricmultisimilaritylossepickitchens100} is the \emph{score-level ensemble of multiple SMS-supervised checkpoints} described in Section~3.3 of that paper but that we did not have the compute budget to reproduce. Concretely, the published recipe combines six independent models: three Adaptive MI-MM variants with different learning-rate schedules and three SMS variants with different relaxation thresholds ($\tau \in \{0.05, 0.10, 0.12\}$). They are all trained for $100$ epochs from the same AVION pretrain~\cite{zhao2023avion} and combines them by directly summing their test similarity matrices, achieving an additional $+1.1$\,mAP and $+0.6$\,nDCG over the best single model. Replicating this ensemble on top of our AVION+SMS pipeline is a natural and well-defined extension due to the fact that row-wise normalization and equal-weight averaging used in Approach~2 (Eq.~\eqref{eq:ensemble}) provide a drop-in replacement for the direct sum, and the graph-diffusion re-ranking of Eq.~\eqref{eq:rerank} can be stacked on top at no training cost. The rationale is supported by three observations from our trajectory and the related literature. First, longer training of a single ViT-L instance yields rapidly diminishing returns past 10--16 epochs, so additional compute is better invested in a few short complementary models than in extending a single one indefinitely. Second, model-soup-style strategies---in which several models fine-tuned with different hyperparameters but the same initialization are combined---consistently match or exceed the best individual model in image classification~\cite{wortsman2022modelsoups} and have already been shown to work in this exact retrieval setting by~\cite{wang2024symmetricmultisimilaritylossepickitchens100}. Third, and our Approach~2 already validates the inference-time fusion machinery on weak frozen baselines, so deploying it on stronger fine-tuned ViT-L+SMS instances is essentially mechanical. The principal obstacle is computational: training six ViT-L+SMS instances for $100$ epochs each requires roughly $40\times$ the GPU-hours of our iteration~6 and is not feasible on a single GPU within a typical project timeline. We therefore frame this direction as the natural continuation of the present work, conditional on multi-GPU infrastructure, in line with the $4\times 4$ RTX~6000~Ada configuration used by Wang et al.\ for ViT-B and the 48\,GB GPUs used for ViT-L. Beyond the ensemble direction, two further extensions follow from the limitations above: integrating the context-refinement mechanism of CR-CLIP~\cite{he2025crclip} on top of our AVION+SMS pipeline (closing a documented portion of the gap to the public leaderboard), and replacing the SMS loss in Eq.~\eqref{eq:sms} with a directly differentiable nDCG surrogate adapted to graded relevance (Smooth-AP~\cite{brown2020smoothap} or a SoDeep-style ranker~\cite{engilberge2019sodeep}).
